\section*{NS-85: full logarithmic taper and the retained prime sum}

This note tests a prescribed full coefficient family, rather than the
projected next-block templates of NS-73. It does not claim a new
approximation theorem. The logarithmic taper is classical; see
Bettin--Conrey--Farmer, arXiv:1211.5191, Theorem 1. Their optimal
unsmoothed asymptotic assumes RH and a separate inverse-zeta-derivative
moment bound. Neither hypothesis is imported into this finite test.

Let $L=\log N$, $N\geq2$, and
\[
 c_n=-\mu(n)\frac{\log(N/n)}{L},\qquad
 R_N=\tau-\sum_{n\leq N}c_n\sigma_n,\qquad
 \eta_N=\sum_{n\leq N}\frac{c_n}{n},\qquad
 S_N=\sum_{n\leq N}\frac{\mu(n)}n\log(N/n)=-L\eta_N.
\]
Here $c_1=-1$ and $c_N=0$; these coefficients are not a fit.
For $1\leq m\leq N$ the complete divisor defect of NS-75 is
\[
 \delta_m=\mathbf1_{m=1}+\sum_{n\mid m}c_n
 =-\frac{\Lambda(m)}L.
\]
This includes $m=1$, with $\Lambda(1)=0$. Indeed
$\sum_{d\mid m}\mu(d)=\mathbf1_{m=1}$ and
$\sum_{d\mid m}\mu(d)\log d=-\Lambda(m)$.
The second identity follows by differentiating
$\sum_{d\mid m}\mu(d)d^z=\prod_{p\mid m}(1-p^z)$ at $z=0$:
one distinct prime gives $-\log p$, and two or more give zero.
Consequently, for every real $1\leq t\leq N$,
\begin{equation}
 L R_N(1/t)=S_Nt-\Psi_1(t),\qquad
 \Psi_1(t)=\sum_{m\leq t}\Lambda(m)\log(t/m)
          =\int_1^t\frac{\psi(u)}u\,du.
\end{equation}
All sums here are finite. A term entering at an integer endpoint has zero
logarithmic factor, so there is no endpoint correction to hide.
For $0<t<1$ the residual is $-\eta_Nt$. Thus its exact partial energy is
\[
 \eta_N^2+\int_1^N |R_N(1/t)|^2\frac{dt}{t^2}
 =\frac1{L^2}\left[S_N^2+
       \int_1^N\left|S_N-\frac{\Psi_1(t)}t\right|^2dt\right].
\]
The remaining interval $t>N$ is part of the complete norm and cannot be
omitted. In particular, a necessary condition for this family's norm to
tend to zero is that the displayed bracket be $o(L^2)$. The ordinary
prime number theorem's pointwise $\Psi_1(t)/t\to1$, by itself, supplies
no such growing-interval mean-square bound and no bound on the remaining
tail. The identity is not the missing arithmetic estimate.

The finite computation sums every exact reciprocal cell through $M$ and
retains the entire tail via NS-78. In its notation,
\[
 u=1-\tfrac12\sum c_n,\quad
 v=\tfrac12\sum c_n\log(n/(2\pi)),\quad
 D=\tfrac16\sum n|c_n|,\quad Q=u\log M+v,
\]
\[
 T=(Q^2+2uQ+2u^2)/M,\quad U=D^2/(3M^3),\qquad
 \left|\int_M^\infty |R_N(1/t)|^2\frac{dt}{t^2}-T\right|
 \leq2\sqrt{TU}+U.
\]
All signed cell cross terms and the main-tail/remainder cross term are
included. Finite coefficient support is not finite physical support.

Two distinct one-coefficient controls are also recorded. Write
$g=\langle R_N,\sigma_1\rangle$, $G_{11}=\|\sigma_1\|^2$, and
$F_N=\|R_N\|^2$. Replacing $c_1$ by $c_1+d$ gives
\[
 \|R_N-d\sigma_1\|^2=F_N-2dg+d^2G_{11}.
\]
The optimal one-atom adjustment is $d=g/G_{11}$, with error
$F_N-g^2/G_{11}$. Exact exterior cancellation uses $d=-\eta_N$ and
has error $F_N+2\eta_Ng+\eta_N^2G_{11}$. These are different rules;
neither is the full $N$-coefficient optimum. Their pairings use the
complete certified arithmetic Gram from NS-61, not a model kernel.

At $N=16,64,256,1024$, both 256- and 384-bit enclosures agree. The
$N=1024$ replay doubles both $M$ and the inherited Gram series cutoff.
The errors decrease on this finite list but substantially exceed the
previously certified optimum at $N=256$. No monotonicity for all $N$,
uniform rate, convergence, or asymptotic failure is inferred. The
exterior-balanced rule also remains inferior to the one-atom optimum
at each tested size; removing one tail term is not a full norm estimate.
